\documentclass[12pt,a4paper]{article}

\usepackage[T2A]{fontenc}
\usepackage[utf8]{inputenc}
\usepackage[russian]{babel}
\usepackage{amsmath,amssymb,amsthm,mathtools}
\usepackage{array}
\usepackage{booktabs}
\usepackage{algorithm}
\usepackage{algpseudocode}
\usepackage{geometry}
\usepackage{enumitem}

\geometry{left=25mm,right=15mm,top=20mm,bottom=20mm}

\newtheorem{definition}{Определение}
\newtheorem{theorem}{Теорема}
\newtheorem{lemma}{Лемма}
\newtheorem{proposition}{Утверждение}
\newtheorem{example}{Пример}
\newtheorem{corollary}{Заклчение}

\DeclareMathOperator{\sgn}{sgn}
\DeclareMathOperator{\seq}{seq}
\DeclareMathOperator{\ord}{ord}

\title{Билет №7 \\ \small Отношения и функции. Отношение эквивалентности и разбиения. Фактор множества. Отношения частичного порядка. (ИТМО)}
\author{Сериков Владислав Александрович}
\date{5 мая 2026}

\begin{document}

\maketitle
\tableofcontents
\newpage

\section{Отношения и функции. Эквивалентности, разбиения, фактор-множества и частичные порядки}

\subsection{Декартово произведение и кортежи}

\begin{definition}
Пусть $A_1,\dots,A_n$ --- множества. \emph{Декартовым произведением}
множеств $A_1,\dots,A_n$ называется множество
\[
A_1\times\cdots\times A_n
=
\{(x_1,\dots,x_n)\mid x_1\in A_1,\dots,x_n\in A_n\}.
\]
Если $A_1=\cdots=A_n=A$, то декартово произведение обозначают через
\[
A^n=\underbrace{A\times\cdots\times A}_{n\ \text{раз}}.
\]
Элементы $(x_1,\dots,x_n)$ называются \emph{упорядоченными $n$-ками}
или \emph{кортежами}.
\end{definition}

Важное отличие кортежа от множества состоит в том, что в кортеже важен порядок координат:
\[
(1,2)\ne (2,1),
\]
хотя
\[
\{1,2\}=\{2,1\}.
\]

\begin{example}
Если $A=\{1,2\}$, $B=\{a,b\}$, то
\[
A\times B=\{(1,a),(1,b),(2,a),(2,b)\},
\]
а
\[
B\times A=\{(a,1),(a,2),(b,1),(b,2)\}.
\]
В общем случае $A\times B\ne B\times A$.
\end{example}

\subsection{Отношения}

\begin{definition}
Пусть $A_1,\dots,A_n$ --- множества. \emph{$n$-местным отношением}
на множествах $A_1,\dots,A_n$ называется любое подмножество
\[
R\subseteq A_1\times\cdots\times A_n.
\]
Если $A_1=\cdots=A_n=A$, то говорят, что $R$ является $n$-местным отношением
на множестве $A$.
\end{definition}

Наиболее важный случай --- бинарные отношения.

\begin{definition}
\emph{Бинарным отношением} между множествами $A$ и $B$ называется любое подмножество
\[
R\subseteq A\times B.
\]
Если $(x,y)\in R$, то пишут
\[
xRy.
\]
Если $A=B$, то говорят, что $R$ является бинарным отношением на $A$.
\end{definition}

\begin{example}
На множестве $\mathbb Z$ можно рассмотреть отношение делимости:
\[
a\mid b \iff \exists k\in\mathbb Z:\ b=ak.
\]
Это бинарное отношение на $\mathbb Z$.
\end{example}

\begin{example}
На множестве $\mathbb R$ отношение
\[
\leqslant=\{(x,y)\in\mathbb R^2\mid x\leq y\}
\]
является бинарным отношением.
\end{example}

\subsection{Область определения, область значений и обратное отношение}

\begin{definition}
Пусть $R\subseteq A\times B$ --- бинарное отношение.

\emph{Областью определения} отношения $R$ называется множество
\[
\delta_R=\{x\in A\mid \exists y\in B:\ (x,y)\in R\}.
\]

\emph{Областью значений} отношения $R$ называется множество
\[
\rho_R=\{y\in B\mid \exists x\in A:\ (x,y)\in R\}.
\]
\end{definition}

\begin{definition}
\emph{Обратным отношением} к отношению $R\subseteq A\times B$ называется отношение
\[
R^{-1}\subseteq B\times A,
\]
заданное правилом
\[
R^{-1}=\{(y,x)\mid (x,y)\in R\}.
\]
\end{definition}

\begin{example}
Пусть
\[
A=\{1,2,3\},\qquad B=\{a,b,c\},
\]
\[
R=\{(1,a),(1,b),(3,c)\}.
\]
Тогда
\[
\delta_R=\{1,3\},
\qquad
\rho_R=\{a,b,c\},
\]
\[
R^{-1}=\{(a,1),(b,1),(c,3)\}.
\]
\end{example}

\subsection{Композиция бинарных отношений}

\begin{definition}
Пусть
\[
R\subseteq A\times B,\qquad S\subseteq B\times C.
\]
\emph{Композиция} отношений $R$ и $S$ --- это отношение
\[
R\circ S\subseteq A\times C,
\]
заданное правилом
\[
R\circ S
=
\{(x,z)\in A\times C\mid \exists y\in B:\ (x,y)\in R,\ (y,z)\in S\}.
\]
\end{definition}

Иногда композицию отношений записывают как $RS$.

\begin{example}
Пусть
\[
R=\{(1,a),(2,a),(2,b)\}\subseteq \{1,2\}\times\{a,b\},
\]
\[
S=\{(a,x),(b,y)\}\subseteq \{a,b\}\times\{x,y\}.
\]
Тогда
\[
R\circ S=\{(1,x),(2,x),(2,y)\}.
\]
Действительно:
\[
1Ra,\ aSx \Rightarrow 1(R\circ S)x,
\]
\[
2Ra,\ aSx \Rightarrow 2(R\circ S)x,
\]
\[
2Rb,\ bSy \Rightarrow 2(R\circ S)y.
\]
\end{example}

\begin{theorem}[Свойства обратного отношения и композиции]
Для любых бинарных отношений, для которых соответствующие композиции определены, выполняются:
\[
(R^{-1})^{-1}=R,
\]
\[
(R\circ S)^{-1}=S^{-1}\circ R^{-1},
\]
\[
(R\circ S)\circ T=R\circ(S\circ T).
\]
\end{theorem}

\begin{proof}
Докажем первое равенство:
\[
(x,y)\in (R^{-1})^{-1}
\iff
(y,x)\in R^{-1}
\iff
(x,y)\in R.
\]
Следовательно,
\[
(R^{-1})^{-1}=R.
\]

Докажем второе равенство. Для произвольных $z$ и $x$ имеем:
\[
(z,x)\in (R\circ S)^{-1}
\iff
(x,z)\in R\circ S.
\]
По определению композиции это эквивалентно существованию такого $y$, что
\[
(x,y)\in R,\qquad (y,z)\in S.
\]
Это, в свою очередь, равносильно
\[
(y,x)\in R^{-1},\qquad (z,y)\in S^{-1}.
\]
Значит,
\[
(z,x)\in S^{-1}\circ R^{-1}.
\]
Следовательно,
\[
(R\circ S)^{-1}=S^{-1}\circ R^{-1}.
\]

Докажем ассоциативность. Пусть
\[
R\subseteq A\times B,\quad S\subseteq B\times C,\quad T\subseteq C\times D.
\]
Тогда
\[
(x,w)\in (R\circ S)\circ T
\]
тогда и только тогда, когда существует $z\in C$ такое, что
\[
(x,z)\in R\circ S,\qquad (z,w)\in T.
\]
Первое условие означает, что существует $y\in B$ такое, что
\[
(x,y)\in R,\qquad (y,z)\in S.
\]
Итак,
\[
(x,w)\in (R\circ S)\circ T
\]
равносильно существованию $y\in B$ и $z\in C$ таких, что
\[
(x,y)\in R,\qquad (y,z)\in S,\qquad (z,w)\in T.
\]
Но это же условие равносильно
\[
(x,w)\in R\circ(S\circ T).
\]
Следовательно,
\[
(R\circ S)\circ T=R\circ(S\circ T).
\]
\end{proof}

\subsection{Матрица бинарного отношения на конечных множествах}

Пусть
\[
A=\{a_1,\dots,a_m\},\qquad B=\{b_1,\dots,b_n\},
\]
и пусть $R\subseteq A\times B$.

\begin{definition}
\emph{Матрицей бинарного отношения} $R$ называется матрица
\[
[R]=(r_{ij})
\]
размера $m\times n$, где
\[
r_{ij}=
\begin{cases}
1, & (a_i,b_j)\in R,\\
0, & (a_i,b_j)\notin R.
\end{cases}
\]
\end{definition}

\begin{example}
Пусть
\[
A=\{1,2,3\},
\]
\[
R=\{(1,1),(1,3),(2,2),(3,1),(3,3)\}\subseteq A^2.
\]
Тогда
\[
[R]=
\begin{pmatrix}
1&0&1\\
0&1&0\\
1&0&1
\end{pmatrix}.
\]
\end{example}

\subsection{Специальные свойства бинарных отношений}

Пусть $R\subseteq A^2$ --- бинарное отношение на множестве $A$.

\begin{definition}
Отношение $R$ называется \emph{рефлексивным}, если
\[
\forall x\in A:\ xRx.
\]
Эквивалентно:
\[
\operatorname{id}_A\subseteq R,
\]
где
\[
\operatorname{id}_A=\{(x,x)\mid x\in A\}.
\]
\end{definition}

\begin{definition}
Отношение $R$ называется \emph{симметричным}, если
\[
\forall x,y\in A:\ xRy\Rightarrow yRx.
\]
Эквивалентно:
\[
R^{-1}=R.
\]
\end{definition}

\begin{definition}
Отношение $R$ называется \emph{антисимметричным}, если
\[
\forall x,y\in A:\ (xRy\ \text{и}\ yRx)\Rightarrow x=y.
\]
\end{definition}

\begin{definition}
Отношение $R$ называется \emph{транзитивным}, если
\[
\forall x,y,z\in A:\ (xRy\ \text{и}\ yRz)\Rightarrow xRz.
\]
Эквивалентно:
\[
R\circ R\subseteq R.
\]
\end{definition}

\begin{example}
На $\mathbb R$ отношение $\leq$ является рефлексивным, антисимметричным и транзитивным, но не является симметричным.
\end{example}

\begin{example}
На любом множестве $A$ отношение равенства $=$ является рефлексивным, симметричным, антисимметричным и транзитивным.
\end{example}

\begin{example}
На множестве людей отношение ``быть знакомым'' обычно симметрично, но не обязательно рефлексивно и не обязательно транзитивно.
\end{example}

\subsection{Функции}

\begin{definition}
Пусть $A$ и $B$ --- множества. \emph{Функцией} или \emph{отображением}
из $A$ в $B$ называется бинарное отношение
\[
f\subseteq A\times B,
\]
такое что:

\begin{enumerate}
\item для любого $x\in A$ существует $y\in B$, для которого $(x,y)\in f$;
\item если $(x,y_1)\in f$ и $(x,y_2)\in f$, то $y_1=y_2$.
\end{enumerate}

Иными словами, каждому элементу $x\in A$ функция ставит в соответствие ровно один элемент $y\in B$.
\end{definition}

Запись:
\[
f:A\to B.
\]
Если $(x,y)\in f$, то пишут
\[
y=f(x).
\]

\begin{definition}
Множество $A$ называется \emph{областью определения} функции $f$, а множество $B$ --- \emph{множеством значений} или \emph{кодоменом}.
Множество
\[
f(A)=\{f(x)\mid x\in A\}
\]
называется \emph{образом} функции $f$.
\end{definition}

Важно различать кодомен $B$ и образ $f(A)$:
\[
f(A)\subseteq B,
\]
но не всегда $f(A)=B$.

\begin{example}
Функция
\[
f:\mathbb R\to\mathbb R,\qquad f(x)=x^2
\]
имеет кодомен $\mathbb R$, но образ
\[
f(\mathbb R)=[0,+\infty).
\]
\end{example}

\subsection{Инъекции, сюръекции и биекции}

\begin{definition}
Функция $f:A\to B$ называется \emph{инъективной} или \emph{инъекцией}, если разные элементы области определения имеют разные образы:
\[
\forall x_1,x_2\in A:\ f(x_1)=f(x_2)\Rightarrow x_1=x_2.
\]
Эквивалентно:
\[
x_1\ne x_2\Rightarrow f(x_1)\ne f(x_2).
\]
\end{definition}

\begin{definition}
Функция $f:A\to B$ называется \emph{сюръективной} или \emph{сюръекцией}, если
\[
f(A)=B.
\]
То есть
\[
\forall y\in B\ \exists x\in A:\ f(x)=y.
\]
\end{definition}

\begin{definition}
Функция $f:A\to B$ называется \emph{биективной} или \emph{биекцией}, если она одновременно инъективна и сюръективна.
\end{definition}

\begin{example}
Функция
\[
f:\mathbb R\to\mathbb R,\qquad f(x)=2x+1
\]
является биекцией.

Действительно, если
\[
2x_1+1=2x_2+1,
\]
то
\[
x_1=x_2,
\]
значит, функция инъективна.

Для любого $y\in\mathbb R$ можно взять
\[
x=\frac{y-1}{2},
\]
и тогда
\[
f(x)=2\cdot \frac{y-1}{2}+1=y.
\]
Значит, функция сюръективна.
\end{example}

\begin{example}
Функция
\[
f:\mathbb R\to\mathbb R,\qquad f(x)=x^2
\]
не является инъективной, так как
\[
f(1)=f(-1)=1,
\]
и не является сюръективной, так как, например, $-1$ не имеет прообраза.
\end{example}

\subsection{Композиция функций}

\begin{definition}
Пусть
\[
f:A\to B,\qquad g:B\to C.
\]
\emph{Композиция} функций $f$ и $g$ --- это функция
\[
g\circ f:A\to C,
\]
заданная правилом
\[
(g\circ f)(x)=g(f(x)).
\]
\end{definition}

В терминах отношений композиция функций является частным случаем композиции бинарных отношений.

\begin{theorem}
Если $f:A\to B$ и $g:B\to C$ --- функции, то $g\circ f:A\to C$ также является функцией.
\end{theorem}

\begin{proof}
Возьмем произвольный $x\in A$.
Так как $f$ --- функция, существует единственный элемент $y\in B$, такой что
\[
y=f(x).
\]
Так как $g$ --- функция, существует единственный элемент $z\in C$, такой что
\[
z=g(y).
\]
Следовательно, каждому $x\in A$ поставлен в соответствие единственный элемент
\[
z=g(f(x)).
\]
Значит, $g\circ f$ является функцией.
\end{proof}

\begin{theorem}
Композиция функций ассоциативна: если
\[
f:A\to B,\qquad g:B\to C,\qquad h:C\to D,
\]
то
\[
h\circ(g\circ f)=(h\circ g)\circ f.
\]
\end{theorem}

\begin{proof}
Для любого $x\in A$ имеем:
\[
(h\circ(g\circ f))(x)=h((g\circ f)(x))=h(g(f(x))).
\]
С другой стороны,
\[
((h\circ g)\circ f)(x)=(h\circ g)(f(x))=h(g(f(x))).
\]
Значения функций совпадают на каждом $x\in A$, следовательно,
\[
h\circ(g\circ f)=(h\circ g)\circ f.
\]
\end{proof}

\subsection{Обратная функция}

\begin{definition}
Пусть $f:A\to B$ --- биекция. Тогда обратное отношение
\[
f^{-1}\subseteq B\times A
\]
является функцией
\[
f^{-1}:B\to A.
\]
Эта функция называется \emph{обратной функцией} к $f$.
\end{definition}

\begin{theorem}
Функция $f:A\to B$ имеет обратную функцию $f^{-1}:B\to A$ тогда и только тогда, когда $f$ является биекцией.
\end{theorem}

\begin{proof}
Докажем необходимость. Пусть $f$ имеет обратную функцию $f^{-1}:B\to A$.
Тогда для каждого $y\in B$ существует $x=f^{-1}(y)$, для которого
\[
f(x)=y.
\]
Следовательно, $f$ сюръективна.

Докажем инъективность. Пусть
\[
f(x_1)=f(x_2).
\]
Обозначим это общее значение через $y$:
\[
y=f(x_1)=f(x_2).
\]
Так как $f^{-1}$ является функцией, у элемента $y$ есть единственный образ:
\[
f^{-1}(y)=x_1
\]
и
\[
f^{-1}(y)=x_2.
\]
Следовательно,
\[
x_1=x_2.
\]
Значит, $f$ инъективна. Итак, $f$ биективна.

Докажем достаточность. Пусть $f$ биективна.
Рассмотрим обратное отношение
\[
f^{-1}=\{(y,x)\mid f(x)=y\}.
\]
Так как $f$ сюръективна, для каждого $y\in B$ найдется $x\in A$, такой что $f(x)=y$.
Так как $f$ инъективна, такой $x$ единственен.
Следовательно, $f^{-1}$ каждому $y\in B$ ставит в соответствие ровно один $x\in A$.
Значит, $f^{-1}:B\to A$ является функцией.
\end{proof}

\subsection{Отношения эквивалентности}

\begin{definition}
Бинарное отношение $E\subseteq A^2$ называется \emph{отношением эквивалентности}, если оно:

\begin{enumerate}
\item рефлексивно:
\[
\forall x\in A:\ xEx;
\]
\item симметрично:
\[
\forall x,y\in A:\ xEy\Rightarrow yEx;
\]
\item транзитивно:
\[
\forall x,y,z\in A:\ (xEy\ \text{и}\ yEz)\Rightarrow xEz.
\]
\end{enumerate}
\end{definition}

Эквивалентность часто обозначают символами
\[
\sim,\quad \equiv,\quad E.
\]

\begin{example}
На любом множестве $A$ отношение равенства является отношением эквивалентности.
\end{example}

\begin{example}
На множестве целых чисел $\mathbb Z$ отношение сравнимости по модулю $n$:
\[
a\equiv b\pmod n
\]
является отношением эквивалентности.

Напомним:
\[
a\equiv b\pmod n \iff n\mid (a-b).
\]
\end{example}

\begin{proof}
Проверим свойства.

Рефлексивность:
\[
a-a=0,
\]
а $n\mid 0$, значит,
\[
a\equiv a\pmod n.
\]

Симметричность:
если
\[
a\equiv b\pmod n,
\]
то
\[
n\mid(a-b).
\]
Но
\[
b-a=-(a-b),
\]
поэтому
\[
n\mid(b-a),
\]
значит,
\[
b\equiv a\pmod n.
\]

Транзитивность:
если
\[
a\equiv b\pmod n,\qquad b\equiv c\pmod n,
\]
то
\[
n\mid(a-b),\qquad n\mid(b-c).
\]
Тогда
\[
a-c=(a-b)+(b-c),
\]
следовательно,
\[
n\mid(a-c).
\]
Значит,
\[
a\equiv c\pmod n.
\]
\end{proof}

\subsection{Классы эквивалентности}

\begin{definition}
Пусть $E$ --- отношение эквивалентности на множестве $A$.
\emph{Классом эквивалентности} элемента $a\in A$ называется множество
\[
[a]_E=\{x\in A\mid xEa\}.
\]
Если отношение понятно из контекста, пишут просто $[a]$.
\end{definition}

\begin{example}
На $\mathbb Z$ рассмотрим сравнимость по модулю $3$.
Тогда существуют три класса эквивалентности:
\[
[0]=\{\dots,-6,-3,0,3,6,\dots\},
\]
\[
[1]=\{\dots,-5,-2,1,4,7,\dots\},
\]
\[
[2]=\{\dots,-4,-1,2,5,8,\dots\}.
\]
\end{example}

\begin{theorem}[Основное свойство классов эквивалентности]
Пусть $E$ --- отношение эквивалентности на $A$.
Тогда для любых $a,b\in A$ либо
\[
[a]_E=[b]_E,
\]
либо
\[
[a]_E\cap[b]_E=\varnothing.
\]
Иначе говоря, любые два класса эквивалентности либо совпадают, либо не пересекаются.
\end{theorem}

\begin{proof}
Пусть
\[
[a]_E\cap[b]_E\ne\varnothing.
\]
Тогда существует элемент $c\in A$, такой что
\[
c\in [a]_E
\quad\text{и}\quad
c\in [b]_E.
\]
Значит,
\[
cEa,\qquad cEb.
\]
Из симметричности получаем
\[
aEc.
\]
Из транзитивности из $aEc$ и $cEb$ следует
\[
aEb.
\]

Докажем, что $[a]_E\subseteq [b]_E$.
Пусть $x\in [a]_E$.
Тогда
\[
xEa.
\]
Так как $aEb$, по транзитивности
\[
xEb.
\]
Значит,
\[
x\in [b]_E.
\]
Следовательно,
\[
[a]_E\subseteq [b]_E.
\]

Аналогично доказывается включение
\[
[b]_E\subseteq [a]_E.
\]
Значит,
\[
[a]_E=[b]_E.
\]

Итак, если классы пересекаются, то они совпадают. Следовательно, если они не совпадают, то они не пересекаются.
\end{proof}

\subsection{Разбиения множества}

\begin{definition}
Пусть $A$ --- непустое множество.
Семейство непустых подмножеств
\[
\mathcal P=\{A_i\mid i\in I\}
\]
называется \emph{разбиением} множества $A$, если:

\begin{enumerate}
\item каждое $A_i$ непусто:
\[
A_i\ne\varnothing;
\]
\item разные элементы семейства не пересекаются:
\[
i\ne j\Rightarrow A_i\cap A_j=\varnothing;
\]
\item объединение всех элементов семейства равно $A$:
\[
A=\bigcup_{i\in I} A_i.
\]
\end{enumerate}
\end{definition}

Иными словами, разбиение --- это представление множества в виде попарно непересекающихся непустых блоков.

\begin{example}
Множество
\[
A=\{1,2,3,4,5,6\}
\]
можно разбить на блоки
\[
\mathcal P=\{\{1,3,5\},\{2,4,6\}\}.
\]
Это разбиение на нечетные и четные элементы.
\end{example}

\subsection{Связь эквивалентностей и разбиений}

\begin{theorem}[Эквивалентности и разбиения]
Пусть $E$ --- отношение эквивалентности на множестве $A$.
Тогда множество всех классов эквивалентности
\[
A/E=\{[a]_E\mid a\in A\}
\]
является разбиением множества $A$.

Обратно, если $\mathcal P=\{A_i\mid i\in I\}$ --- разбиение множества $A$, то отношение $E_{\mathcal P}$, заданное правилом
\[
xE_{\mathcal P}y
\iff
\exists i\in I:\ x\in A_i\ \text{и}\ y\in A_i,
\]
является отношением эквивалентности на $A$.

Таким образом, отношения эквивалентности на $A$ взаимно соответствуют разбиениям множества $A$.
\end{theorem}

\begin{proof}
Сначала пусть $E$ --- отношение эквивалентности.

Докажем, что $A/E$ является разбиением.

Во-первых, каждый класс непуст:
\[
a\in [a]_E
\]
по рефлексивности, значит,
\[
[a]_E\ne\varnothing.
\]

Во-вторых, каждый элемент $a\in A$ принадлежит некоторому классу, а именно $[a]_E$.
Следовательно,
\[
A=\bigcup_{a\in A}[a]_E.
\]

В-третьих, по предыдущей теореме любые два класса эквивалентности либо совпадают, либо не пересекаются.
Следовательно, различные классы попарно не пересекаются.

Значит, $A/E$ --- разбиение $A$.

Теперь пусть $\mathcal P=\{A_i\mid i\in I\}$ --- разбиение $A$.
Зададим отношение:
\[
xE_{\mathcal P}y
\iff
x \text{ и } y \text{ лежат в одном блоке разбиения}.
\]

Проверим свойства эквивалентности.

Рефлексивность.
Так как $\mathcal P$ покрывает $A$, для любого $x\in A$ существует $i\in I$, такое что
\[
x\in A_i.
\]
Следовательно, $x$ лежит в одном блоке с самим собой:
\[
xE_{\mathcal P}x.
\]

Симметричность.
Если
\[
xE_{\mathcal P}y,
\]
то существуют $i\in I$, такое что
\[
x,y\in A_i.
\]
Тогда также
\[
y,x\in A_i,
\]
значит,
\[
yE_{\mathcal P}x.
\]

Транзитивность.
Пусть
\[
xE_{\mathcal P}y,\qquad yE_{\mathcal P}z.
\]
Тогда существуют $i,j\in I$, такие что
\[
x,y\in A_i,
\]
\[
y,z\in A_j.
\]
Так как $y\in A_i\cap A_j$, а блоки разбиения попарно не пересекаются, получаем
\[
A_i=A_j.
\]
Следовательно,
\[
x,z\in A_i,
\]
значит,
\[
xE_{\mathcal P}z.
\]

Таким образом, $E_{\mathcal P}$ --- отношение эквивалентности.
\end{proof}

\subsection{Фактор-множество}

\begin{definition}
Пусть $E$ --- отношение эквивалентности на множестве $A$.
\emph{Фактор-множеством} множества $A$ по отношению $E$ называется множество всех классов эквивалентности:
\[
A/E=\{[a]_E\mid a\in A\}.
\]
\end{definition}

\begin{example}
Для отношения сравнимости по модулю $n$ на $\mathbb Z$ фактор-множество имеет вид
\[
\mathbb Z/n\mathbb Z=\{[0],[1],\dots,[n-1]\}.
\]
\end{example}

\begin{example}
Пусть $A$ --- множество студентов университета, а отношение $E$ задано так:
\[
xEy
\iff
x \text{ и } y \text{ учатся в одной группе}.
\]
Тогда классы эквивалентности --- это учебные группы, а фактор-множество $A/E$ --- множество всех групп.
\end{example}

\subsection{Каноническая проекция на фактор-множество}

\begin{definition}
Пусть $E$ --- отношение эквивалентности на $A$.
Отображение
\[
\pi:A\to A/E,
\qquad
\pi(a)=[a]_E,
\]
называется \emph{канонической проекцией} или \emph{естественным отображением} на фактор-множество.
\end{definition}

\begin{theorem}
Каноническая проекция
\[
\pi:A\to A/E
\]
является сюръекцией.
\end{theorem}

\begin{proof}
Пусть $C\in A/E$.
По определению фактор-множества существует $a\in A$, такое что
\[
C=[a]_E.
\]
Но
\[
\pi(a)=[a]_E=C.
\]
Следовательно, каждый элемент фактор-множества имеет прообраз.
Значит, $\pi$ сюръективна.
\end{proof}

\subsection{Отношения порядка}

Отношения порядка формализуют идею сравнения элементов.

\begin{definition}
Бинарное отношение $\preceq$ на множестве $A$ называется \emph{предпорядком} или \emph{квазипорядком}, если оно рефлексивно и транзитивно.
\end{definition}

\begin{definition}
Бинарное отношение $\leq$ на множестве $A$ называется \emph{частичным порядком}, если оно:

\begin{enumerate}
\item рефлексивно:
\[
\forall x\in A:\ x\leq x;
\]
\item антисимметрично:
\[
(x\leq y\ \text{и}\ y\leq x)\Rightarrow x=y;
\]
\item транзитивно:
\[
(x\leq y\ \text{и}\ y\leq z)\Rightarrow x\leq z.
\]
\end{enumerate}
\end{definition}

Пара
\[
(A,\leq)
\]
называется \emph{частично упорядоченным множеством}.

\begin{example}
Множество $\mathbb R$ с обычным отношением $\leq$ является частично упорядоченным множеством.
\end{example}

\begin{example}
Для любого множества $U$ булеан
\[
\mathcal P(U)
\]
с отношением включения $\subseteq$ является частично упорядоченным множеством.
\end{example}

\begin{proof}
Проверим свойства включения.

Рефлексивность:
\[
A\subseteq A
\]
для любого $A\subseteq U$.

Антисимметричность:
если
\[
A\subseteq B
\quad\text{и}\quad
B\subseteq A,
\]
то
\[
A=B.
\]

Транзитивность:
если
\[
A\subseteq B
\quad\text{и}\quad
B\subseteq C,
\]
то
\[
A\subseteq C.
\]
Следовательно, $\subseteq$ --- частичный порядок.
\end{proof}

\subsection{Сравнимые и несравнимые элементы}

\begin{definition}
Пусть $(A,\leq)$ --- частично упорядоченное множество.
Элементы $x,y\in A$ называются \emph{сравнимыми}, если
\[
x\leq y
\quad\text{или}\quad
y\leq x.
\]
Если ни одно из этих соотношений не выполняется, то элементы называются \emph{несравнимыми}.
\end{definition}

\begin{example}
В множестве $\mathcal P(\{1,2,3\})$ элементы
\[
\{1\},\qquad \{2\}
\]
несравнимы относительно включения:
\[
\{1\}\nsubseteq \{2\},
\qquad
\{2\}\nsubseteq \{1\}.
\]
\end{example}

\subsection{Линейный порядок}

\begin{definition}
Частичный порядок $\leq$ на множестве $A$ называется \emph{линейным порядком}, если любые два элемента сравнимы:
\[
\forall x,y\in A:\ x\leq y\ \text{или}\ y\leq x.
\]
Пара $(A,\leq)$ называется \emph{линейно упорядоченным множеством}.
\end{definition}

\begin{example}
Множества
\[
(\mathbb Z,\leq),\quad (\mathbb Q,\leq),\quad (\mathbb R,\leq)
\]
с обычным порядком являются линейно упорядоченными.
\end{example}

\begin{example}
Множество $\mathcal P(\{1,2,3\})$ с отношением $\subseteq$ не является линейно упорядоченным, поскольку, например, $\{1\}$ и $\{2\}$ несравнимы.
\end{example}

\subsection{Строгий порядок}

\begin{definition}
Пусть $\leq$ --- частичный порядок.
Соответствующий \emph{строгий порядок} $<$ определяется правилом
\[
x<y\iff x\leq y\ \text{и}\ x\ne y.
\]
\end{definition}

Строгий порядок обычно не является рефлексивным:
\[
x<x
\]
никогда не выполняется.

\begin{theorem}
Если $\leq$ --- частичный порядок на $A$, то соответствующее отношение $<$ транзитивно и антирефлексивно.
\end{theorem}

\begin{proof}
Антирефлексивность очевидна:
\[
x<x
\iff
x\leq x\ \text{и}\ x\ne x,
\]
что невозможно.

Докажем транзитивность. Пусть
\[
x<y,\qquad y<z.
\]
Тогда
\[
x\leq y,\quad x\ne y,
\]
и
\[
y\leq z,\quad y\ne z.
\]
По транзитивности $\leq$ получаем
\[
x\leq z.
\]
Остается доказать, что $x\ne z$.
Если бы $x=z$, то из $y\leq z=x$ и $x\leq y$ по антисимметричности следовало бы
\[
x=y,
\]
что противоречит $x\ne y$.
Значит,
\[
x\ne z.
\]
Следовательно,
\[
x<z.
\]
\end{proof}

\subsection{Минимальные, максимальные, наименьшие и наибольшие элементы}

Пусть $(A,\leq)$ --- частично упорядоченное множество.

\begin{definition}
Элемент $a\in A$ называется \emph{минимальным}, если не существует элемента $x\in A$, такого что
\[
x<a.
\]
Эквивалентно:
\[
x\leq a\Rightarrow x=a.
\]
\end{definition}

\begin{definition}
Элемент $a\in A$ называется \emph{максимальным}, если не существует элемента $x\in A$, такого что
\[
a<x.
\]
Эквивалентно:
\[
a\leq x\Rightarrow x=a.
\]
\end{definition}

\begin{definition}
Элемент $a\in A$ называется \emph{наименьшим}, если
\[
\forall x\in A:\ a\leq x.
\]
\end{definition}

\begin{definition}
Элемент $a\in A$ называется \emph{наибольшим}, если
\[
\forall x\in A:\ x\leq a.
\]
\end{definition}

\begin{theorem}
Если в частично упорядоченном множестве существует наименьший элемент, то он единственен.
Аналогично, если существует наибольший элемент, то он единственен.
\end{theorem}

\begin{proof}
Пусть $a$ и $b$ --- два наименьших элемента.
Так как $a$ наименьший, то
\[
a\leq b.
\]
Так как $b$ наименьший, то
\[
b\leq a.
\]
По антисимметричности получаем
\[
a=b.
\]
Значит, наименьший элемент единственен.

Доказательство для наибольшего элемента аналогично.
Если $a$ и $b$ --- наибольшие элементы, то
\[
a\leq b
\]
и
\[
b\leq a,
\]
откуда
\[
a=b.
\]
\end{proof}

\begin{theorem}
Всякий наименьший элемент является минимальным.
Всякий наибольший элемент является максимальным.
\end{theorem}

\begin{proof}
Пусть $a$ --- наименьший элемент.
Тогда
\[
a\leq x
\]
для любого $x\in A$.

Предположим, что $x\leq a$.
Так как $a\leq x$, по антисимметричности получаем
\[
x=a.
\]
Следовательно, не существует элемента, строго меньшего $a$.
Значит, $a$ минимален.

Доказательство для наибольшего элемента аналогично.
\end{proof}

Обратное неверно: минимальный элемент не обязательно является наименьшим.

\begin{example}
Пусть
\[
A=\{\{1\},\{2\},\{1,2\}\}
\]
с порядком $\subseteq$.
Тогда $\{1\}$ и $\{2\}$ --- минимальные элементы.
Но наименьшего элемента нет, поскольку
\[
\{1\}\nsubseteq \{2\}
\]
и
\[
\{2\}\nsubseteq \{1\}.
\]
\end{example}

\subsection{Верхние и нижние грани. Супремум и инфимум}

Пусть $(A,\leq)$ --- частично упорядоченное множество, $B\subseteq A$.

\begin{definition}
Элемент $u\in A$ называется \emph{верхней гранью} множества $B$, если
\[
\forall b\in B:\ b\leq u.
\]
\end{definition}

\begin{definition}
Элемент $l\in A$ называется \emph{нижней гранью} множества $B$, если
\[
\forall b\in B:\ l\leq b.
\]
\end{definition}

\begin{definition}
\emph{Точной верхней гранью} или \emph{супремумом} множества $B$ называется наименьшая из верхних граней.
Обозначение:
\[
\sup B.
\]
\end{definition}

\begin{definition}
\emph{Точной нижней гранью} или \emph{инфимумом} множества $B$ называется наибольшая из нижних граней.
Обозначение:
\[
\inf B.
\]
\end{definition}

\begin{example}
В множестве $\mathbb R$ с обычным порядком для
\[
B=(0,1)
\]
имеем
\[
\inf B=0,
\qquad
\sup B=1.
\]
При этом $0\notin B$ и $1\notin B$.
\end{example}

\subsection{Диаграммы Хассе}

Для конечных частично упорядоченных множеств удобно использовать диаграммы Хассе.

\begin{definition}
Пусть $(A,\leq)$ --- конечное частично упорядоченное множество.
Говорят, что элемент $y$ \emph{покрывает} элемент $x$, если
\[
x<y
\]
и не существует $z\in A$, такого что
\[
x<z<y.
\]
\end{definition}

\emph{Диаграмма Хассе} строится так:
\begin{enumerate}
\item элементы множества изображаются точками;
\item если $y$ покрывает $x$, то $x$ и $y$ соединяют отрезком;
\item точку $y$ располагают выше точки $x$;
\item ребра, следующие из транзитивности, обычно не рисуют.
\end{enumerate}

\begin{example}
Пусть
\[
A=\mathcal P(\{1,2\})
=
\{\varnothing,\{1\},\{2\},\{1,2\}\}
\]
с отношением $\subseteq$.
Диаграмма Хассе имеет вид:
\[
\begin{array}{c}
\{1,2\}\\
/\quad \backslash\\
\{1\}\quad \{2\}\\
\backslash\quad /\\
\varnothing
\end{array}
\]
\end{example}

\subsection{Алгоритм построения классов эквивалентности по матрице отношения}

Пусть $A=\{a_1,\dots,a_n\}$, и отношение эквивалентности $E$ задано матрицей
\[
[E]=(e_{ij}).
\]

\subsubsection*{Шаги алгоритма}

\begin{enumerate}
\item Выбрать первый еще не использованный элемент $a_i$.
\item Найти все индексы $j$, для которых
\[
e_{ij}=1.
\]
\item Сформировать класс
\[
[a_i]_E=\{a_j\mid e_{ij}=1\}.
\]
\item Пометить все элементы найденного класса как использованные.
\item Повторять шаги $1$--$4$, пока не будут использованы все элементы множества $A$.
\end{enumerate}

\begin{example}
Пусть
\[
A=\{1,2,3,4,5\},
\]
и отношение эквивалентности задано матрицей
\[
[E]=
\begin{pmatrix}
1&0&1&0&1\\
0&1&0&1&0\\
1&0&1&0&1\\
0&1&0&1&0\\
1&0&1&0&1
\end{pmatrix}.
\]

Берем элемент $1$.
В первой строке единицы стоят в столбцах $1,3,5$.
Значит,
\[
[1]=\{1,3,5\}.
\]

Следующий неиспользованный элемент --- $2$.
Во второй строке единицы стоят в столбцах $2,4$.
Значит,
\[
[2]=\{2,4\}.
\]

Итак,
\[
A/E=\{\{1,3,5\},\{2,4\}\}.
\]
\end{example}

\subsection{Алгоритм проверки, является ли отношение частичным порядком}

Пусть $A=\{a_1,\dots,a_n\}$, и отношение $R\subseteq A^2$ задано матрицей
\[
[R]=(r_{ij}).
\]

\subsubsection*{Шаги алгоритма}

\begin{enumerate}
\item Проверить рефлексивность:
\[
r_{ii}=1
\]
для всех $i=1,\dots,n$.

\item Проверить антисимметричность:
для всех $i\ne j$ не должно быть одновременно
\[
r_{ij}=1
\quad\text{и}\quad
r_{ji}=1.
\]

\item Проверить транзитивность:
если
\[
r_{ij}=1
\quad\text{и}\quad
r_{jk}=1,
\]
то должно выполняться
\[
r_{ik}=1.
\]

Эквивалентно можно проверить матричное условие:
\[
[R]\cdot [R]\leq [R],
\]
где умножение понимается булево:
\[
1+1=1.
\]

\item Если все три условия выполнены, то $R$ --- частичный порядок.
Если хотя бы одно условие нарушено, то $R$ не является частичным порядком.
\end{enumerate}

\begin{example}
Пусть
\[
A=\{1,2,3\},
\]
и
\[
R=\{(1,1),(2,2),(3,3),(1,2),(1,3),(2,3)\}.
\]
Матрица отношения:
\[
[R]=
\begin{pmatrix}
1&1&1\\
0&1&1\\
0&0&1
\end{pmatrix}.
\]

Рефлексивность есть, так как на диагонали стоят единицы.

Антисимметричность есть, так как вне диагонали нет симметричных пар единиц.

Транзитивность:
\[
1R2,\quad 2R3 \Rightarrow 1R3,
\]
и это действительно есть.
Остальные случаи также выполняются.

Следовательно, $R$ является частичным порядком.
Это обычный порядок
\[
1\leq 2\leq 3.
\]
\end{example}

\subsection{Типичные примеры отношений}

\subsubsection*{1. Равенство}

На любом множестве $A$ отношение
\[
=
\]
является одновременно:
\begin{itemize}
\item эквивалентностью;
\item частичным порядком.
\end{itemize}

\subsubsection*{2. Включение}

На $\mathcal P(U)$ отношение $\subseteq$ является частичным порядком, но обычно не является линейным порядком.

\subsubsection*{3. Сравнение по модулю}

На $\mathbb Z$ отношение
\[
a\equiv b\pmod n
\]
является эквивалентностью, но не является частичным порядком при $n\geq 2$, поскольку не выполняется антисимметричность:
например,
\[
1\equiv 3\pmod 2,
\]
но
\[
1\ne 3.
\]

\subsubsection*{4. Делимость}

На множестве натуральных чисел $\mathbb N$ отношение
\[
a\mid b
\]
является частичным порядком, если считать натуральные числа положительными.

\begin{proof}
Рефлексивность:
\[
a\mid a,
\]
так как
\[
a=1\cdot a.
\]

Антисимметричность:
если
\[
a\mid b
\quad\text{и}\quad
b\mid a,
\]
то существуют $m,n\in\mathbb N$, такие что
\[
b=am,\qquad a=bn.
\]
Тогда
\[
a=amn.
\]
Так как $a>0$, получаем
\[
mn=1.
\]
Следовательно,
\[
m=n=1,
\]
и
\[
a=b.
\]

Транзитивность:
если
\[
a\mid b
\quad\text{и}\quad
b\mid c,
\]
то
\[
b=ak,\qquad c=bl.
\]
Тогда
\[
c=akl,
\]
значит,
\[
a\mid c.
\]
\end{proof}

\subsection{Краткое сравнение эквивалентностей и порядков}

\[
\begin{array}{c|c|c}
\text{Свойство} & \text{Эквивалентность} & \text{Частичный порядок}\\
\hline
\text{Рефлексивность} & + & +\\
\text{Симметричность} & + & \text{обычно нет}\\
\text{Антисимметричность} & \text{обычно нет} & +\\
\text{Транзитивность} & + & +
\end{array}
\]

Отношение эквивалентности группирует элементы в классы, считая элементы одного класса ``одинаковыми'' относительно выбранного признака.

Отношение порядка организует элементы по принципу сравнения, но не обязательно позволяет сравнить любую пару элементов.

\subsection{Минимальный набор фактов для запоминания}

\begin{enumerate}
\item Бинарное отношение между $A$ и $B$ --- это подмножество $A\times B$.
\item Функция $f:A\to B$ --- это специальное бинарное отношение, в котором каждому $x\in A$ соответствует ровно один $y\in B$.
\item Эквивалентность --- это рефлексивное, симметричное и транзитивное отношение.
\item Классы эквивалентности либо совпадают, либо не пересекаются.
\item Эквивалентности на $A$ взаимно соответствуют разбиениям $A$.
\item Фактор-множество $A/E$ --- это множество всех классов эквивалентности.
\item Частичный порядок --- это рефлексивное, антисимметричное и транзитивное отношение.
\item Линейный порядок --- это частичный порядок, в котором любые два элемента сравнимы.
\item Наименьший элемент, если существует, единственен и является минимальным.
\item Наибольший элемент, если существует, единственен и является максимальным.
\end{enumerate}

\end{document}
